\section{光のフラッシュで放出するドーパミン}
\label{sec:dishdopamine}

培養にドーパミンを教師として与えた直接の実験として私たちが知っている唯一のものは、研究者が遠隔で接続するオルガノイドのプラットフォームで行われた\cite{Jordan2024}。オルガノイドを浸す液に、光感受性の「ケージ」に閉じ込めたドーパミンを加えた。ケージに入った分子は受容体に作用しない。ケージドドーパミンの濃度は $1$~mM である。ネットワークに報酬を与えたいときは紫外線を当てる。ケージが壊れ、ドーパミンが空間に出ていく。

ループは次のようにできている。同じ刺激を何度も与え、それが以前より多くのスパイクを引き起こしたら、フラッシュを点けてドーパミンを放出する。$240$ 回の繰り返しで、誘発スパイクの数は全体として増えた。著者ら自身、この実験一つだけでは増加がドーパミンによるとは結論できないと書いている\cite{Jordan2024}。

エンジニアにとって、これはシャーレの中に三要素則\eqref{eq:threefactor}を組む最初の試みである。送り手と受け手の活動が適格度トレースを残し、共通の信号が変化を定着させるかどうかを決める。しかしここでの信号は正にしかならない。報酬があるかないかであり、負の誤差 $\delta<0$ をこの装置は出せない。そのうえドーパミンは広がり、ゆっくりとしか除去されない。\eqref{eq:lowpass}の濃度と同じである。そのため頻繁なフラッシュは単に全体のレベルを上げるだけかもしれない。学習をこれと区別するには二つの対照が必要である。ネットワークの応答と無関係なランダムな時刻に同じ数のフラッシュを与える対照と、ケージドドーパミンなしで同じフラッシュを与える対照である。さらに休止後の確認も必要である。ドーパミンが去ったあとも変化が残っているかどうかである。

\section{ドーパミンがシリコンを教える}
\label{sec:biohybrid}

逆向きの実験では、生きた細胞が電子回路を教える\cite{Keene2020}。ドーパミンを分泌する細胞を、有機ニューロモルフィック素子の上の微小流路で育てる。この素子はトランジスタで、そのチャネルのコンダクタンスがシナプスの重みの役をする。細胞から出たドーパミンは素子の電極で酸化され、それがコンダクタンスを変える。この装置はシナプス間隙から伝達物質を除去する仕組みも再現しているので、重みは長く変えておくことも、元に戻すこともできる。

回路としては、これは化学入力の漏れ積分器である。重みはドーパミンの流入を蓄積し、「汲み出し」が忘れる速さを決める。\eqref{eq:lowpass}と同じ RC 回路である。ここで肝心なのは接続の向きである。第\ref{sec:debugging}章のプローブにブロードキャストレジスタが見えないのは、それが化学的だからである。この素子はまさにその化学レジスタを読み取り、電気的な重みに変える。ブレイン・コンピュータ・インターフェースにとって、これはデータバスだけでなく制御レジスタも聴き取れるセンサーへの道である。

\section{自前のドーパミンニューロンをもつオルガノイド}
\label{sec:midbrainorg}

ヒトの幹細胞から、\nt{DOPA} ニューロンのある脳領域に似たオルガノイドを育てることができる。そこには働くドーパミンニューロンがあり、ドーパミンを分泌する\cite{Jo2016}。こうしたオルガノイドは主に病気の研究のために育てられている。そのドーパミンを別のネットワークの教師として使った実験は、私たちには見つからなかった。

生きたニューロンでできたマシンにとって、これは足りない部品である。第\ref{sec:livingmachine}章のテストベンチは、制御レジスタなしのデータバスとフロー制御だった。ここにはブロードキャスト信号の源がすでにある。足りないのは批評家、つまり予測誤差\eqref{eq:ctd}を計算してこれらのニューロンを制御するネットワークである。大脳皮質のオルガノイドをこうしたオルガノイドとつなぎ、誤差に応じてドーパミンが放出されるようにすることは未解決の課題であり、今のところ実験ではなく仮説である。

\section{ドーパミンなしで棒を立てるオルガノイド}
\label{sec:cartpole}

最近の実験で最も完全なものは、ドーパミンをまったく使わない\cite{Robbins2026}。マウス大脳皮質のオルガノイドを、「台車の上の棒」（倒立振子）の課題と閉ループでつないだ。オルガノイドの活動が台車を動かし、棒の位置は刺激として戻される。試行の合間に、オルガノイドは短い高頻度の学習刺激を受ける。それをどれにするかは強化学習プログラムが選ぶ。

結果は測定されている。
\begin{itemize}
\item 大多数のオルガノイドで、プログラムが選んだ学習信号は、ランダムに選んだ信号や信号なしより良い成績をもたらした。
\item $45$ 分の休止後、改善は残らなかった。
\item AMPA と NMDA の両方のグルタミン酸受容体を遮断すると、改善はなくなった。
\end{itemize}

最後の項目は、学習したものがどこにあるかを語る。\nt{GLUT} シナプスの中であり、\ref{sec:weightstore}節で入力と出力の同時検出器として働く受容体を通じてである。2番目の項目は何が足りないかを示す。速い変化はあるが定着がない。第\ref{sec:motorlearning}章の、遅い学習者を欠いた速い学習者と同じである。そして教師の役割は変わった。電流パターンを選ぶエンジニアはプログラムに置き換わったが、教師は依然としてシャーレの外にいる。

電極アレイ上の培養を学習させたもっと古い実験の中にも、ドーパミンを使ったものは一つも見つからなかった。学習信号はどれも電気刺激だった。

\section{誰が培養を教えるのか}

\begin{table}[t]
\centering
\footnotesize
\setlength{\tabcolsep}{4pt}
\renewcommand{\arraystretch}{1.15}
\begin{tabular}{>{\raggedright}p{2.7cm}>{\raggedright}p{2.5cm}>{\raggedright}p{3.2cm}>{\raggedright\arraybackslash}p{4.4cm}}
\toprule
実験 & 教えるのは誰か & 学習信号 & わかっていること \\
\midrule
目的の応答で刺激を止める & エンジニア & 刺激の終了 & 応答が定着する（測定済み） \\
DishBrain（第\ref{sec:dishbrain}章） & エンジニア & 予測可能な電流またはランダム電流 & セッション内のラリーの有意な伸びは完全なフィードバックのときだけ、無音では効果が小さい（測定済み）。セッション間では不安定。原因は議論中 \\
台車の上の棒 & 強化学習プログラム & プログラムが選ぶ高頻度刺激 & ランダムより良いが、休止後に残らない（測定済み） \\
フラッシュで放出するドーパミン & 「スパイクが増えたら報酬」という規則 & 光で解放されるドーパミン & スパイクの増加は観察。ドーパミンの役割は未証明 \\
バイオハイブリッドシナプス & 生きた細胞 & 細胞からのドーパミン & 素子の重みの長期変化と復帰（測定済み） \\
\nt{DOPA} ニューロンをもつオルガノイド & --- & --- & ドーパミン源はある。教師としては未使用 \\
\bottomrule
\end{tabular}
\caption{チップ上の生きたニューロンを誰が何で教えるか。}
\label{tab:dishteachers}
\end{table}

培養自体が学習するどの実験でも、教師は今のところ外にいる（表\ref{tab:dishteachers}）。エンジニアか、プログラムか、エンジニアが定めた規則である。ドーパミンがシャーレに入ったのは一度だけで、その役割はまだ証明されていない。第\ref{sec:dofa}章のように批評家、誤差、適格度トレースを備えた本物のブロードキャストチャネルをシャーレの中に組むことが次の一歩であり、まだ誰も踏み出していない。
